\section{Incidence and the quaternionic Pontryagin class}
\label{inc:section}

The lower-volume comparison and its metric equality case are now
proved. To obtain the quaternionic upper bound, we first determine
the characteristic class of an actual bundle supplied by the cut
geometry: the normal four-plane bundle of cut incidence.
The later loop-energy construction will use this same bundle,
not only its sphere-fibration homotopy type.
We use the standard integral bundle and cohomology tools in
\citet{Hatcher2002} and \citet{MilnorStasheff1974}; the geometric
constructions and the relations needed here are proved below.
In this section and the next three, suppose that $M$ has quaternionic
type, so $n=4r$, $r\ge2$, and $k=3$. Choose an integral generator
$u\in H^4(M;\Z)$ and an orientation with $\int_Mu^r=1$.
Write $C_p=\operatorname{Cut}(p)$. The cell attachment in
Section~\ref{geom:section} gives
\[
 H^*(C_p;\Z)=\Z[u|_{C_p}]/((u|_{C_p})^r).
\]
Indeed restriction is an isomorphism in every degree below $4r-1$,
and $C_p$ has dimension $4r-4$. In particular it is orientable,
and its fundamental class is dual to $\pm u$ in $M$.

\subsection{The actual normal bundle}

\begin{proposition}\label{inc:normal}
The cut incidence
\[
 F=\{(p,q)\in M\times M:q\in C_p\}
\]
is a smooth embedded submanifold. Both projections are smooth
fiber bundles with fiber a cut locus. If $u_1,u_2$ are the two
pullbacks of $u$, its rank-four normal bundle $\mathcal V$ can be oriented
so that
\begin{equation}\label{inc:normal-classes}
 e(\mathcal V)=u_1-u_2,\qquad p_1(\mathcal V)=2s(u_1+u_2),\qquad s\in2\Z+1.
\end{equation}
Projection of $\mathcal V_{(p,q)}$ into $T_pM$ identifies it with
$N_pC_q$. The unit sphere of this four-plane is exactly the
great sphere of initial directions of minimizing geodesics
from $p$ to $q$.
\end{proposition}

\begin{proof}
The full cut map
\[
 c:SM\longrightarrow M\times M,\qquad
 c(\xi)=(\pi(\xi),\exp_{\pi(\xi)}(D\xi))
\]
has constant rank $4r+(4r-4)=8r-4$. Its fibers are the connected
great three-spheres from Proposition~\ref{geom:structure}.
Here constant rank and compactness give an embedded image,
as follows. Constant-rank charts give local image germs through
any image point. Charts overlapping along its fiber have the
same germ, so connectedness of that fiber makes all these germs
agree. Cover the compact fiber by finitely many such charts.
Properness excludes additional inverse images outside their
union over a sufficiently small neighborhood of the image point.
Thus the whole image in that neighborhood is the one smooth
submanifold germ. This proves the assertion. The first projection
is a submersion because its composition with $c$ is $\pi$;
symmetry gives the same conclusion for the second. Proper
submersions are smooth fiber bundles.

For $(a,b)\in \mathcal V_{(p,q)}\subset T_pM\oplus T_qM$, orthogonality
to $(T_pC_q,0)$ gives $a\in N_pC_q$. If $a=0$, surjectivity
of the second projection on $T_{(p,q)}F$ forces $b=0$.
The projection to $N_pC_q$ is therefore injective and, both
spaces having dimension four, an isomorphism. First variation
for the constant distance $D$ shows that the initial directions
from $p$ to $q$ are normal to $C_q$. Their great sphere has
the full dimension of $S(N_pC_q)$, proving the last assertion.
In particular, if $q_1,q_2\in C_p$ are distinct, then
\begin{equation}\label{inc:disjoint-planes}
 N_pC_{q_1}\cap N_pC_{q_2}=\{0\}:
\end{equation}
a common nonzero direction would reach both points at time $D$.

The cut loci and $M$ are simply connected, so the fibration
makes $F$ simply connected. Its normal bundle is orientable.
Leray--Hirsch \citep[Theorem~4D.1]{Hatcher2002}, using
$1,u_2,\ldots,u_2^{r-1}$ on the fibers,
and disjointness from the diagonal give the integral presentation
\begin{equation}\label{inc:ring}
 H^*(F;\Z)=\Z[u_1,u_2]/(u_1^{r+1},\Delta(u_1,u_2)),
 \qquad \Delta(a,b)=\sum_{j=0}^r a^jb^{r-j}.
\end{equation}
To justify completeness, the displayed relation is the restriction
of the diagonal class in $M\times M$. It is monic of degree $r$
in $u_2$, so the proposed quotient has exactly the Leray--Hirsch
basis. The dual of $F$ in $M\times M$ restricts on either factor
slice to $\pm u$. Its restriction to the disjoint diagonal is
zero. Its two factor coefficients are consequently opposite,
and after an orientation choice its restriction to $F$ is $u_1-u_2$.
The self-intersection formula gives the asserted Euler class.

The interchange of factors preserves the unoriented normal
bundle. Since $u_1,u_2$ are an integral basis of $H^4(F;\Z)$,
we have $p_1(\mathcal V)=t(u_1+u_2)$ for an integer $t$.
Also $H^2(F;\Z/2)=0$, so $\mathcal V$ is spin.
We use the classical identification
$\operatorname{Spin}(4)=\operatorname{SU}(2)\times\operatorname{SU}(2)$
and its two complex two-dimensional half-spin representations;
see \citet[Sections~2.1--2.2]{Moore2010}. For an oriented spin
four-plane bundle, its two complex half-spin bundles have
second Chern classes
\[
 \frac{-p_1(\mathcal V)+2e(\mathcal V)}4,\qquad
 \frac{-p_1(\mathcal V)-2e(\mathcal V)}4.
\]
This follows directly from the formal half-spin weights
$\pm(t_1+t_2)/2$ and $\pm(t_1-t_2)/2$, where
$p_1=t_1^2+t_2^2$ and $e=t_1t_2$.
Integrality in the basis $u_1,u_2$ implies $t\equiv2\pmod4$.
Writing $t=2s$ proves the proposition.
\end{proof}

Changing $u$ to $-u$ replaces $s$ by $-s$. Reverse the orientation
of $\mathcal V$ at the same time to retain
$e(\mathcal V)=u_1-u_2$, and reverse the orientation of $M$
exactly when $r$ is odd to retain $\int_Mu^r=1$.
We may therefore take $s\ge1$.

\subsection{The rank constraint from two direction planes}

For distinct $q_1,q_2\in C_p$, \eqref{inc:disjoint-planes} places
two four-planes with zero intersection inside $T_pM$. Their direct
sum has a complement of rank $4(r-2)$, which limits the degree
of its total Pontryagin class. We will use this rank restriction
to determine $p(TM)$ from $s$. The parameter space must therefore
exclude coincident cut points; put
\[
 Z=\{(p,q_1,q_2):q_1,q_2\in C_p,\ q_1\ne q_2\}.
\]
Use $v,x_1,x_2$ for the pullbacks of $u$ from its three points.
In this section polynomial degree is one quarter of cohomological
degree. In degrees divisible by four the rational cohomology of
$Z$ is the algebra
\begin{equation}\label{inc:triple-ring}
 A=\mathbb Q[v,x_1,x_2]/(v^{r+1},\Delta(v,x_1),S(v,x_1,x_2)),
 \quad
 S(v,x_1,x_2)=\sum_{i+j+\ell=r-1}v^ix_1^jx_2^\ell.
\end{equation}
Here and below this assertion identifies the subalgebra consisting
of these degrees; no claim about other degrees is needed.

For completeness, before removing $q_1=q_2$, Leray--Hirsch gives
the relations $v^{r+1}$, $\Delta(v,x_1)$ and $\Delta(v,x_2)$.
The relative diagonal has codimension $4(r-1)$. Its dual is
$\pm S$. To see this, the dual is a homogeneous polynomial of
degree $r-1$ annihilated by $x_1-x_2$. In degree $r$ the only
relations are the two displayed $\Delta$ polynomials; their
difference is $(x_1-x_2)S$. More explicitly, multiplying a
representative of the dual by $x_1-x_2$ gives a constant linear
combination of those two relations. Setting $x_1=x_2$ makes
the two constants opposite. Cancellation of $x_1-x_2$ then
takes place in the ordinary polynomial ring, before passage
to cohomology, and shows that the dual is proportional to $S$.
Restriction to a fixed-$p$ fiber gives the ordinary diagonal
class of $C_p\times C_p$, fixing the factor as $\pm1$.
The Thom sequence for removing the relative diagonal is
surjective onto the complement in degrees divisible by four,
because the relevant next cohomology group of $F$ is zero.
Restriction from the ambient fiber product to that diagonal is
surjective, so the projection formula identifies the kernel
with the ideal of its dual. The relation $\Delta(v,x_2)$ then
becomes redundant, proving \eqref{inc:triple-ring}.
Successive division by its two monic relations proves that $A$
is free over $\mathbb Q[v]/(v^{r+1})$, with basis
\begin{equation}\label{inc:free-basis}
 \{x_1^a x_2^b:0\le a<r,\ 0\le b<r-1\}.
\end{equation}

\begin{theorem}\label{inc:pontryagin}
The total Pontryagin class of $M$ is
\begin{equation}\label{inc:pontryagin-formula}
 p(TM)=\frac{(1+2su+u^2)^{r+1}}{1+4su}
 \quad\text{in }H^*(M;\mathbb Q).
\end{equation}
The quotient means its power-series expansion truncated by
$u^{r+1}=0$.
\end{theorem}

\begin{proof}
Over $Z$, the two actual normal four-planes in
\eqref{inc:disjoint-planes} have direct sum inside the pullback
of $TM$ from $p$. Choose a complementary bundle $W$ of rank
$4(r-2)$. Set
\[
 D_s(v,x)=1+2s(v+x)+(v-x)^2,
 \qquad p_0(v)=\frac{(1+2sv+v^2)^{r+1}}{1+4sv}.
\]
The rank-four identity $p_2=e^2$
\citep[Corollary~15.8]{MilnorStasheff1974} and
Proposition~\ref{inc:normal} give
\begin{equation}\label{inc:complement}
 p(W)=\frac{p(TM)(v)}{D_s(v,x_1)D_s(v,x_2)}.
\end{equation}
Its polynomial degree is at most $t_0=2r-4$.

We will prove that the formal characteristic class
\[
 \mathcal W_0=\frac{p_0(v)}{D_s(v,x_1)D_s(v,x_2)}\in A
\]
has degree at most $t_0$ and a degree-$t_0$ component that is nonzero
modulo $v$. The freeness in
\eqref{inc:free-basis} will then show that any first discrepancy
between $p(TM)(v)$ and $p_0(v)$ violates the rank bound on $W$.
For the degree bound we use only the algebra
\eqref{inc:triple-ring}, without presuming the answer for $TM$.
Successive monic polynomial divisions factor
\[
 T^{r+1}=(T-v)(T-x_1)(T-x_2)
                   \prod_{j=4}^{r+1}(T-z_j)
\]
in an extension of $A$ obtained by adjoining roots of monic
polynomials. Give each root degree one; the monic relations
are homogeneous for this grading. Each adjoining algebra is
free with a basis containing $1$, so this extension is injective. The positive-degree
elementary symmetric functions of all $r+1$ roots vanish.
Consequently, the product of $D_s(v,z)$ over all these roots is
$(1+2sv+v^2)^{r+1}$. The factor at $z=v$ is $1+4sv$.
Dividing these constant-one series therefore gives
\[
 \mathcal W_0=\prod_{j=4}^{r+1}D_s(v,z_j).
\]
There are $r-2$ factors, each of degree at most two, as required.

We also need the top component $(\mathcal W_0)_{t_0}$ to be nonzero
modulo $v$. For $r=2$ it is $1$. For $r\ge3$, let
\[
 (1+2sx+x^2)^{-1}=\sum_{j\ge0}a_jx^j.
\]
The recursion $a_0=1$, $a_1=-2s$ and
$a_j=-2sa_{j-1}-a_{j-2}$ gives
\begin{equation}\label{inc:reciprocal-identity}
 a_j^2-a_{j-1}a_{j+1}=1\qquad(j\ge1).
\end{equation}
Indeed the left side is unchanged on replacing $j$ by $j-1$,
and its initial value is one. Modulo $v$ the algebra can also be
presented as
\[
 \mathbb Q[x_1,x_2]/
 \left(x_1^r,x_2^r,\sum_{j=0}^{r-1}x_1^jx_2^{r-1-j}\right).
\]
In degree $2r-4$ the surviving monomials after the first two
truncations are
\[
 x_1^{r-3}x_2^{r-1},\qquad
 x_1^{r-2}x_2^{r-2},\qquad x_1^{r-1}x_2^{r-3}.
\]
The remaining ideal gives exactly their sum as a relation:
multiplying its generator by any monomial of degree $r-3$
gives this same sum. Their coefficients in $\mathcal W_0$ are, in order,
$a_{r-3}a_{r-1},a_{r-2}^2,a_{r-3}a_{r-1}$.
Equation~\eqref{inc:reciprocal-identity} makes the middle
coefficient differ from either outside coefficient by one.
Thus the component modulo $v$ is the nonzero class of
$x_1^{r-2}x_2^{r-2}$, one of the monomials in
\eqref{inc:free-basis}.

Finally suppose that $p(TM)(v)/p_0(v)$ first differs from $1$
in degree $i$, where $1\le i\le r$. Write its first coefficient
as $b\ne0$. In degree $t_0+i$, modulo $v^{i+1}$,
\eqref{inc:complement} is $b v^i(\mathcal W_0)_{t_0}$.
This is nonzero by \eqref{inc:free-basis} and the preceding
calculation: in a free module over $\mathbb Q[v]/(v^{r+1})$,
a coefficient nonzero modulo $v$ remains nonzero after
multiplication by $v^i$, modulo $v^{i+1}$, for $i\le r$.
It contradicts the rank bound on $W$.
Therefore there is no such $i$, which proves the theorem.
\end{proof}
